\section{Boolean rigidity at every admissible size}\label{sec:rigidity}

We specialize the exact configuration count to sign labels and integer
local equations. Fix an alternative sign eigenvector $\tau$. It is balanced
by \eqref{eq:planted}, so it differs from $\sigma$ on the same number
$k$ of vertices in each planted class. Replacing $\tau$ by $-\tau$
allows $0\le k\le N/2$. Set $\alpha=k/N$. The root laws are
\begin{equation}\label{eq:sign-laws}
 \mu_+(1)=1-\alpha,\qquad \mu_-(1)=\alpha,
 \qquad \Var(\mu_\pm)=4\alpha(1-\alpha).
\end{equation}
We treat macroscopic $k$ by a uniform entropy deficit and small $k$
by a support exploration.

\subsection{The two-star exponent and its strict variance deficit}
Use the exact ordered-star pairing count of
\eqref{eq:fmp-exact-types}, now on two old types with $a$ internal
and $b$ cross slots, and restrict to the integer sign equation.
Let $P_+,P_-$ be the two profile laws and let $Q_+,Q_-$ be the
internal pair laws, with cross pair law $R$. Their exponent per $N$ is
\[
 \Phi=H(P_+)+H(P_-)-\frac a2[H(Q_+)+H(Q_-)]-bH(R).
\]
Symmetrizing internal slots and cross slots separately preserves the
pooled tables and the legal equation, and increases this exponent.
For these symmetric laws omit one internal neighbor, which is uniquely
determined by the root and all retained neighbors. Let $J_+,J_-$
be their conditional total correlations, and write
\[
 I_+=2H(\mu_+)-H(Q_+),\quad I_-=2H(\mu_-)-H(Q_-),\quad
 I_\times=H(\mu_+)+H(\mu_-)-H(R).
\]
The entropy chain rule, using the appropriate root marginal at each
neighbor, gives exactly
\begin{equation}\label{eq:entropy-identity}
 -\Phi=\left(\frac a2-1\right)(I_++I_-)
                      +bI_\times+J_++J_-.
\end{equation}
The coefficient of $I_++I_-$ in \eqref{eq:entropy-identity} is zero
at $a=2$. It is therefore useful to retain the internal conditional
law in an auxiliary product distribution instead of forcing it to be
independent of the root.

\begin{lemma}[Quantitative sign deficit]\label{lem:sign-deficit}
For every feasible sign type with overlap $\alpha$ as in
\eqref{eq:sign-laws}, every $a>b>0$ satisfies
\begin{equation}\label{eq:sign-deficit}
 \Phi\le-\frac{2b^2}{a^4}\alpha^2(1-\alpha)^2.
\end{equation}
\end{lemma}
\begin{proof}
By the preceding symmetrization it suffices to bound the larger
exponent of the symmetric profile laws. Let $P$ be the
law of the retained plus star
$(X,Y_1,\ldots,Y_{a-1},Z_1,\ldots,Z_b)$. Define $T$ as follows.
Keep the same root law. Given that root, draw the $a-1$ internal
neighbors independently with their original conditional law
$\mathcal L(Y_1\mid X)$. Draw the $b$ cross neighbors independently
with the root-independent law $\mu_-$, independently of the internal
neighbors. Expanding relative entropy gives
\begin{equation}\label{eq:KL-star}
 \KL(P\Vert T)=J_++bI_\times\le L,
 \qquad L=-\Phi,
\end{equation}
by \eqref{eq:entropy-identity} and $a\ge2$.
Every positive $P$ atom has positive corresponding conditional and
marginal probabilities under $T$, so this identity is valid with zero
marginals as well. Conditional laws at zero-probability roots can be
chosen arbitrarily.

Put
\[
 m(X)=\E[Y_1\mid X],\quad v(X)=\Var(Y_1\mid X),\quad
 U=\delta X-\sum_{i<a}Y_i-\sum_{j=1}^bZ_j.
\]
Under $P$, the integer constraint makes $U$ the omitted internal sign
$Y_a$. Exchangeability of the internal slots therefore gives
$\E_P(U-m(X))^2=\E v(X)$. Under $T$, conditional independence gives
\[
 \Var_T(U\mid X)=(a-1)v(X)+b\Var(\mu_-).
\]
The squared conditional-mean contribution is nonnegative. Consequently
\begin{equation}\label{eq:variance-difference}
 \E_T(U-m(X))^2-\E_P(U-m(X))^2
 \ge(a-2)\E v(X)+b\Var(\mu_-)
 \ge b\Var(\mu_-).
\end{equation}
On both distributions $|U|\le\delta+(a-1)+b=2a-1$ and $|m(X)|\le1$.
The function $(U-m(X))^2$ therefore lies in $[0,4a^2]$.
The bounded-range expectation inequality and Pinsker's inequality yield
\[
 b\Var(\mu_-)
 \le4a^2\|P-T\|_{\TV}
 \le4a^2\sqrt{\KL(P\Vert T)/2}
 \le4a^2\sqrt{L/2}.
\]
Substituting \eqref{eq:sign-laws} and squaring proves
\eqref{eq:sign-deficit}. The proof did not require positive internal
mutual-information coefficients and thus includes $(a,b)=(2,1)$.
\end{proof}

Let $Z_k$ count sign eigenvectors with $k$ flips in each block.
Restrict the exact ordered-profile count to those root counts. The root
laws are then exactly \eqref{eq:sign-laws}. For this fixed finite sign
alphabet and fixed degrees there are polynomially many star types.
The exact pairing/bijection factorial weights have uniform
$O_{a,b}(\log(N+1))$ remainders, including zero entries, as in
Section~\ref{sec:frame-moment-proof}. Lemma~\ref{lem:sign-deficit}
therefore gives, for fixed constants $C_0,C_1$,
\begin{equation}\label{eq:Boolean-fiber}
 \E Z_k\le C_0(N+1)^{C_1}
 \exp\!\left[-\frac{2b^2}{a^4}N\alpha^2(1-\alpha)^2\right].
\end{equation}
This is an integer-sign count; it does not reduce the two signs
modulo two. All choices of root signs are already in the ordered-star multinomials. For any fixed $\eta>0$, summing over the macroscopic range gives
\begin{equation}\label{eq:macro}
 \E\#\{\tau:A\tau=\delta\tau,\ \eta N\le k\le N/2\}
 \le e^{-c_{a,b,\eta}N}
\end{equation}
for sufficiently large $N$. The strictly negative exponential dominates
the polynomial prefactor and the $N$ possible overlaps. Conditioning
on simplicity costs only a bounded factor.

\subsection{Uniform small-support exclusion}
\begin{lemma}[Small-support eigenvectors]\label{lem:small-support}
For fixed $a>b>0$ there is $\eps=\eps(a,b)>0$ such that, in the
uniform-simple model,
\begin{equation}\label{eq:small-support}
 \Prob\{\exists v\ne0:Av=\delta v,\ |\supp v|\le\eps N\}
 =O_{a,b}\bigl(N^{-(D-2)(\delta+1)}\bigr).
\end{equation}
\end{lemma}
\begin{proof}
We first obtain deterministic necessary conditions in a simple
$D$-regular graph. Since $0<\delta<D$, every such eigenvector has
coordinate sum zero and has both positive and negative entries.
At its largest positive entry, the eigenvector equation requires at
least $\delta$ positive neighbors: each is at most that entry and all
other contributions are nonpositive. The analogous observation at its
most negative entry requires at least $\delta$ negative neighbors.
Simplicity makes these neighbors distinct from their roots. Thus
\begin{equation}\label{eq:min-support}
 |\supp v|\ge q_0:=2(\delta+1).
\end{equation}
Let $S=\supp v$, $|S|=q$, and let $T$ be its exterior neighbor set.
Every vertex of $T$ has at least two incident edges from $S$; one
nonzero summand cannot satisfy its zero-coordinate eigenvector equation.

For a fixed $S$, explore the configuration pairings incident to its
$Dq$ stubs, skipping any stub whose partner has already been exposed.
Call an exposure a collision if its partner lies in $S$ or at an
exterior vertex already reached. If $i$ exposed edges lie in $S$,
$j$ leave $S$, and $t$ distinct exterior vertices are reached, then
\[
 Dq=2i+j,\qquad t\le j/2,\qquad
 \#\text{collisions}=i+j-t\ge Dq/2.
\]
This counts an edge within $S$ only once, though it consumes two
support stubs.

For $q\le\eps N$ and sufficiently small fixed $\eps$, each relevant
internal-pairing or cross-bijection pool retains $\Theta_{a,b}(N)$
stubs throughout this exploration. There are only $O_D(q)$ stubs at
vertices of $S$ and previously reached exterior vertices. Conditional
on every previous exposure, each collision has probability at most
$C_{a,b}q/N$. The exposure rule can choose any unpaired support stub;
uniform pairings and bijections retain the required conditional law.
Pad the exploration to $Dq$ steps with deterministic noncollisions.
A union bound over sets of at least $Dq/2$ collision positions gives
\[
 \Prob\{S\text{ has the necessary exterior condition}\}
 \le2^{Dq}(C_{a,b}q/N)^{Dq/2}.
\]
Multiplying by $\binom{2N}{q}$ and enlarging a fixed constant bounds
the probability of any such set of size $q$ by
\begin{equation}\label{eq:collision-bound}
 f_q:=\left[K_{a,b}\left(\frac qN\right)^h\right]^q,
 \qquad h=(D-2)/2>0.
\end{equation}
Indeed
$f_{q+1}/f_q\le K_{a,b}e^h((q+1)/N)^h$.
Choose $\eps$ so this is at most $1/2$ throughout the range.
Summing from $q_0$ to $\eps N$ gives
$O(N^{-hq_0})=O(N^{-(D-2)(\delta+1)})$.

The minimum-support condition \eqref{eq:min-support} was used only
for the desired simple-graph event. The necessary-set exploration
was bounded in the unconditioned configuration model. Conditioning
its three factors on simplicity costs a bounded factor and gives
exactly the requested law, proving the lemma.
\end{proof}

\begin{proof}[Proof of Theorem~\ref{thm:boolean}]
For an alternative with $1\le k\le\eps N/2$, the nonzero eigenvector
$\tau-\sigma$ has support $2k\le\eps N$ and is excluded by
Lemma~\ref{lem:small-support}. The remaining alternatives satisfy
$\eps N/2\le k\le N/2$ and are excluded with exponentially small
probability by \eqref{eq:macro}. Combining the two bounds proves
\eqref{eq:boolean-rate}. The smallest possible exponent is two,
attained at $(a,b)=(2,1)$.
\end{proof}
